\subsection{Prediction Visualization}
\label{sec:visualization}

We visualize the spatial prediction quality to complement the aggregate metrics in Table~\ref{tab:main_results}. TSC-Fusion predictions are evaluated against observations for both temperature and salinity across multiple depth levels.

Fig.~\ref{fig:prediction_maps} shows mean-field predictions for TEMP and SALT averaged over all test samples at three representative depths (surface, 50~m, and 200~m). The model reproduces the large-scale spatial structure---the meridional temperature gradient, the coastal salinity patterns, and the subsurface thermocline features---with high fidelity. Residual errors are predominantly localized near frontal boundaries and coastal regions where sharp gradients challenge any data-driven method. The spatial error distribution is consistent across depths, indicating that the multi-branch fusion preserves vertical structure.

\begin{figure*}[t]
\centering
\includegraphics[width=0.94\textwidth]{figures/fig_prediction_maps.pdf}
\caption{Prediction maps of TSC-Fusion averaged over the test set at three depth levels. Left columns: ground truth and prediction for TEMP; right columns: same for SALT. The model captures large-scale spatial structure with errors concentrated near frontal and coastal regions.}
\label{fig:prediction_maps}
\end{figure*}

To assess vertical structure preservation, we compute depth-wise RMSE profiles for TSC-Fusion, FuXi-Ocean, and AxiomOcean over all test samples. Fig.~\ref{fig:depth_rmse} (left) shows that TSC-Fusion maintains lower RMSE across the entire water column for temperature, with the largest relative improvement in the thermocline region (50--300~m) where vertical gradients are strongest. For salinity (right), the advantage is concentrated in the upper 200~m where water-mass variability is highest, consistent with the TSC memory module's design.

The T-S scatter diagram in Fig.~\ref{fig:depth_rmse} (center) further demonstrates that TSC-Fusion predictions preserve the joint temperature-salinity relationship. The predicted T-S distribution closely follows the observed climatological curve, indicating physically consistent coupling between the two variables---a property that is not directly enforced by the training objective but emerges from the thermohaline-structured architecture.

\begin{figure*}[t]
\centering
\includegraphics[width=0.94\textwidth]{figures/fig_depth_rmse.pdf}
\caption{Left: Depth-wise RMSE profiles for TEMP and SALT comparing TSC-Fusion with FuXi-Ocean and AxiomOcean. TSC-Fusion maintains lower error throughout the water column, with the largest advantage in the thermocline region. Right: T-S scatter diagram comparing observed and predicted temperature-salinity pairs. The predicted distribution closely follows the climatological relationship.}
\label{fig:depth_rmse}
\end{figure*}
